\section{Related Work}
\label{sec:related}

Vegas sits at the intersection of three lines of work --- domain
specific languages for smart contracts, automated game-theoretic
analysis, and event-structure semantics for concurrent processes.
None of these alone covers Vegas's design, and the value of the
language design is in their combination.

\subsection{Smart-contract domain-specific languages}

Several existing DSLs target contract correctness.
Marlowe~\citep{Lamela2020Marlowe} offers a small financial
contract calculus for Cardano with formal semantics in
Isabelle/HOL; its scope is multi-party financial agreements and
its proofs are about safety and conservation, not about strategic
properties of the players.
Scilla~\citep{Sergey2019Scilla} provides a typed intermediate
language for Zilliqa with Coq-verified semantics and a strict
separation of communication and computation; like Marlowe, its
guarantees concern functional behavior of the contract rather
than the game it implements.
Move~\citep{blackshear2019move} introduces linear resource types
to make ``how much can leave the contract'' a property of the
type system.

Vegas differs in two structural ways.  First, the source
language describes a \emph{game played by the contract},
not just a contract, and the compiler emits a game-theoretic
artifact (a Gambit \TT{.efg}) alongside the executable code.
Second, the language is consciously narrower than these
alternatives: no recursion, no token transfers, no arbitrary
external calls.  This narrowness is what gives the multi-backend
story its consistency.

Verification tools that check smart contracts against temporal
or safety specifications --- VerX~\citep{Permenev2020Verx},
KEVM~\citep{park2018formal}, Certora~\citep{CertoraWhitepaper2022}
--- are complementary.  Their question is ``does this bytecode
satisfy this specification''; Vegas's question is ``which
specification?''  Strategically-aware specifications are exactly
what Vegas's Gambit/MAID/SMT artifacts are.

\subsection{Game-theoretic analysis of smart contracts}

The MEV literature~\citep{Daian2020Flash,qin2022quantifying}
documents extensive strategic activity on existing chains:
front-running, sandwich attacks, transaction reordering for
profit.  These analyses are case studies of deployed contracts,
performed after the fact and largely by hand.  Mitigation work
has typically taken two routes: implementation-level cryptographic
defenses for commit-reveal patterns~\citep{breidenbach2018enter}
and runtime alternatives such as private transaction
pools~\citep{Weintraub2022FlashbotsAuction}.  Vegas's
contribution is to build a specific subclass --- observation-before
action on concurrent public moves --- into the language: the
mitigation is a structural property of the generated contract,
not a discipline applied by humans.

\subsection{Multi-agent influence diagrams and game representations}

MAID~\citep{koller2003multi} provides a factored representation
of multi-agent decisions as a Bayesian network with decision and
utility nodes.  We discussed the MAID backend in
\S\ref{sec:other-backends}; the relevant limit is that plain
MAIDs cannot encode context-dependent action sets, which makes
them an awkward target for Vegas games whose guards are
non-trivial.  Constrained-MAID extensions exist in the literature
but are not standardized.

The Gambit project~\citep{McKelveyMcLennanTurocy2006Gambit} and
OpenSpiel~\citep{LanctotEtAl2019OpenSpiel} provide tools for
analyzing extensive-form games once one is given.  Vegas adds
the path from a programmer's specification to such a game.

\subsection{Event structures and dependency graphs}

The EventGraph is a particular flavour of event
structure~\citep{NielsenPlotkinWinskel1981,Winskel1986EventStructures}
specialized to the protocol setting --- typed, finite,
visibility-aware, with explicit commit/reveal node kinds.  The
program-dependence graph~\citep{ferrante1987program} and
static-single-assignment form~\citep{cytron1991efficiently} are
the immediate technical antecedents on the program-analysis side:
EventGraph edges are computed by the same dependency analysis
those works pioneered, only specialized to the
\TT{commit}/\TT{reveal} discipline.

\subsection{Compositional game theory and IF logic}

Compositional Game Theory~\citep{GhaniHedgesWinschel2018CGT} and
the work it builds on develop a categorical view of open games
that compose like programs.  IF logic~\citep{HintikkaSandu1989,
MannSanduSevenster2011IFL} provides a logical apparatus for
imperfect-information games in which informational independence
is a first-class construct.  Vegas's per-program Diamond DAG
encoding (\S\ref{sec:other:diamond}) has a sympathetic flavour
--- a player's choice schedule is a function on prior witnesses,
which is the IF-logic notion of a Skolem function for an
existential under independence --- but the connection is at the
level of stylistic resonance, not a formal correspondence we
have spelled out.  The companion VegasCore
(\S\ref{sec:vegascore}) does this work on the side of native
kernel-game semantics rather than open-game composition.

\subsection{Multiparty session types}

Multiparty session types~\citep{HondaYoshidaCarbone2008MPST}
specify the legal sequencing of messages in a multi-party
protocol; Castellan and Yoshida's recent work
\citep{CastellanYoshidaSessionGame} relates session types to
game semantics directly.  Vegas's Scribble backend uses session
types as a target, projecting the EventGraph's public-message
subset onto a global type.  The relationship is more
``information set $\subseteq$ session type than the reverse'':
session types record which messages are legal in a given
prefix, but they do not directly record information sets.  A
fuller comparison is future work.

\subsection{Refinement and operational correctness}

The bail-out subgame and the Solidity scheduling layer share
the shape of an implementation refining a more abstract
specification, in the sense of Lynch and
Vaandrager~\citep{LynchVaandrager1995}.  Vegas's companion
formalization VegasCore states a stochastic-step refinement
theorem on the abstract event-graph machine.  A formal proof
that Vegas's Solidity backend refines its EventGraph is on the
project's TODO list; this is precisely the gap the fellowship
proposal addresses.
